\documentclass[11pt]{article}
\usepackage{amsmath, amssymb}
\usepackage[margin=1in]{geometry}
\usepackage{booktabs}
\usepackage{siunitx}

\title{Governing Equations for a Waste-Heat-Driven, Two-Bed SAWH System (No HTF Loop)}
\author{}
\date{}

\begin{document}
\maketitle

\section*{System picture and notation}

Two adsorbent-coated contactors, $A$ and $B$, alternate roles every half-cycle (duration set by desorber RH threshold, capped at $\tau_{1/2,\max}$). During a half-cycle, contactor~$A$ adsorbs from ambient process air while contactor~$B$ desorbs under partial vacuum. This is the same two-bed cycling logic as the sibling \texttt{waste\_heat\_cycle\_lumped} package, but the pumped heat-transfer-fluid (HTF) loop that used to shuttle the heat of adsorption from~$A$ to the desorbing bed~$B$ has been removed. The desorbing bed now couples directly to the fixed waste-heat stream through a single equivalent UA (the series combination of the waste-heat-stream HX and the contactor-side UA that used to sandwich the loop on either end); the adsorbing bed rejects its heat of adsorption to ambient through a finned heat sink instead of having it extracted by the loop. This generalizes Wilson et al.'s single-bed solar model: Eqs.~1--2 and 5--6 carry over in spirit, but the solar/absorber/glass stack is replaced by the direct waste-heat coupling, and a second bed enables continuous operation.

\textbf{Subscripts:} $a$ = adsorbing contactor; $d$ = desorbing contactor (identities swap each half-cycle); $wh$ = waste-heat stream (fixed inlet temperature and mass flow, coupled directly to $d$); $cond$ = condenser; $amb$ = ambient process airstream.

\section{State variables}

Per unit footprint area (\si{\metre\squared}), unless noted.

\textbf{Hydrogel sorbent (default):}
\begin{equation}
\mathbf{y} = \bigl[c_{w,a},\, H_a,\, c_{w,d},\, H_d,\, T_a,\, T_d,\, T_{cond}\bigr]
\end{equation}

\textbf{MOF sorbent (optional):}
\begin{equation}
\mathbf{y} = \bigl[q_a,\, q_d,\, T_a,\, T_d,\, T_{cond}\bigr]
\end{equation}
where $q$ is gravimetric loading (\si{\kilogram\water\per\kilogram\adsorbent}).

There is no loop-fluid state $T_f$: with the pumped HTF loop removed, the only thermal states are the two contactor temperatures and the condenser temperature.

\section{Contactor energy balances}

\textbf{Desorbing contactor} (sealed, partial vacuum $P_{cond} \approx 20$--$60$~mbar), coupled directly to the waste-heat stream:
\begin{equation}
(\rho c_p V)_{d}\,\frac{dT_d}{dt}
= \dot{Q}_{wh\to d}
- \dot{m}_{des}\,h_{des}
- h_{gap}\,A_d\,(T_d - T_{cond})
- \varepsilon_d\sigma A_d\,(T_d^4 - T_{cond}^4)
\label{eq:desorbing}
\end{equation}

\textbf{Adsorbing contactor} (open to process air, near-ambient pressure), rejecting the heat of adsorption to ambient through a finned heat sink rather than to a loop:
\begin{equation}
(\rho c_p V)_{a}\,\frac{dT_a}{dt}
= \dot{m}_{ads}\,h_{ads}
- h_{fin}\,A_a\,(T_a - T_{amb})
\label{eq:adsorbing}
\end{equation}
\begin{equation}
h_{fin} = \phi_{fin}\,h_{amb}
\end{equation}
The adsorbing-bed fin area ratio $\phi_{fin}$ is sized like the condenser's existing fin stack: plain ambient convection alone ($h_{amb}$ only) gives a cooldown time constant of $(\rho c_p V)_a / (h_{amb} A_a) \approx \SI{167}{\minute}$, far slower than a half-cycle, so without fins the bed cannot shed the heat of adsorption plus its residual heat from the prior desorbing role before it must swap roles again. Finning to $\phi_{fin} = 7.1$ cuts that time constant to $\approx \SI{23}{\minute}$.

At partial vacuum, $h_{gap}$ uses a rarefied-gas conductance (Knudsen--continuum blend), not Hollands free convection.

\section{Direct waste-heat coupling}

The waste-heat stream couples to the desorbing bed through a single equivalent UA, $UA_{wh,d}$, formed by combining the waste-heat-stream HX conductance and the contactor-side conductance in series (the two conductances that used to sandwich the now-removed loop fluid):
\begin{equation}
UA_{wh,d} = \left(\frac{1}{UA_{wh}} + \frac{1}{UA_{d}}\right)^{-1} = \frac{UA_{wh}\,UA_d}{UA_{wh} + UA_d}
\label{eq:uaeq}
\end{equation}

NTU--$\varepsilon$ closure ($\dot{m}_{wh}$ is a fixed waste-heat mass flow, not a control variable):
\begin{equation}
\dot{Q}_{wh\to d} = \dot{m}_{wh}\,c_{p,wh}\,(T_{wh,in} - T_d)\,\Bigl[1-\exp\!\Bigl(-\frac{UA_{wh,d}}{\dot{m}_{wh}\,c_{p,wh}}\Bigr)\Bigr]
\label{eq:Qwhd}
\end{equation}

Unlike the looped variant, there is no intermediate loop-fluid temperature and no pumped flow to modulate the coupling: $T_{wh,in}$ and $\dot{m}_{wh}$ are fixed boundary conditions, and $UA_{wh,d}$ is a fixed device parameter.

\section{Condenser energy balance}

\begin{equation}
(\rho c_p L)_{cond}\,\frac{dT_{cond}}{dt}
= h_{gap}\,A_d\,(T_d - T_{cond})
+ \dot{m}_{des}\,h_{fg}
+ \varepsilon_d\sigma\,(T_d^4 - T_{cond}^4)
- A_r\,h_{amb}\,(T_{cond} - T_{amb})
\label{eq:condenser}
\end{equation}

\section{Mass transport}

Unchanged from \texttt{waste\_heat\_cycle\_lumped}: only the thermal coupling to the desorbing bed changes, not the sorbent mass-transport model.

\subsection{Hydrogel (PAM--LiCl, Wilson Eqs.~5--6)}

\textbf{Adsorption (open bed):}
\begin{equation}
\frac{dc_{w,a}}{dt} = \frac{g_{chamber}}{H_0}\,\frac{P_{sat}(T_a)}{RT_a}\,\bigl(RH_{amb} - a_w(c_{w,a}, H_a)\bigr)
\label{eq:massads}
\end{equation}
\begin{equation}
\frac{dH_a}{dt} = g_{chamber}\,\frac{MW_w}{\rho_{sol}}\,\frac{P_{sat}(T_a)}{RT_a}\,\bigl(RH_{amb} - a_w\bigr)
\end{equation}
\begin{equation}
\dot{m}_{ads} = MW_w\,\bigl(\dot{c}_{w,a}\,H_a + c_{w,a}\,\dot{H}_a\bigr), \qquad \dot{m}_{ads} \ge 0
\end{equation}
Wilson's printed Eq.~6 omits the mass-transfer coefficient, but $\frac{MW_w}{\rho_{sol}}\frac{P_{sat}}{RT}$ alone is dimensionless, so $dH/dt$ as printed has no units of m/s. The $g_{chamber}$ (or $g_{gap}$, desorption) factor is required for dimensional consistency and to match the desorption mass-rate identity $\dot{m}_{des} = -\rho_{sol}A_c\,dH/dt$ used elsewhere in Wilson et al.; this implementation includes it.

\textbf{Desorption (vacuum-limited Wilson inventory):}

Vacuum conductance sets the primary desorption capacity:
\begin{equation}
\dot{m}_{vac} = C_{vac}(t)\,\bigl[P_{sat}(T_d) - P_{cond}\bigr],
\qquad
\dot{m}_{vac} \le \frac{MW_w\,c_{w,d}\,H_d}{\tau_{dep,\max}}
\label{eq:massdes}
\end{equation}
where $\tau_{dep,\max}$ caps depletion when gel inventory is low.

Wilson heat--mass-analogy rates (same Eq.~5--6 structure as absorption, with vapor-gap $g$ from Hollands natural convection between gel and condenser):
\begin{equation}
c_r(T_d, T_{cond}) = \frac{P_{sat}(T_{cond})}{P_{sat}(T_d)}\,\frac{T_d}{T_{cond}}
\label{eq:crdes}
\end{equation}
\begin{equation}
\left.\frac{dc_{w,d}}{dt}\right|_{\mathrm{Wilson}}
= \frac{g_{gap}(H_d)}{H_0}\,\frac{P_{sat}(T_d)}{RT_d}\,
\Bigl[c_r(T_d, T_{cond}) - a_w(c_{w,d}, H_d)\Bigr]
\label{eq:massdes_cw}
\end{equation}
\begin{equation}
\left.\frac{dH_d}{dt}\right|_{\mathrm{Wilson}}
= g_{gap}(H_d)\,\frac{MW_w}{\rho_{sol}}\,\frac{P_{sat}(T_d)}{RT_d}\,
\Bigl[c_r(T_d, T_{cond}) - a_w(c_{w,d}, H_d)\Bigr]
\label{eq:massdes_H}
\end{equation}
\begin{equation}
\dot{m}_{Wilson} = MW_w\,\Bigl(
\left.\frac{dc_{w,d}}{dt}\right|_{\mathrm{Wilson}} H_d
+ c_{w,d}\,\left.\frac{dH_d}{dt}\right|_{\mathrm{Wilson}}\Bigr)
\label{eq:mwilson}
\end{equation}

The gel inventory rates applied in the ODE are scaled to the vacuum limit:
\begin{equation}
\lambda = \min\!\Bigl(1,\; \frac{\dot{m}_{vac}}{\dot{m}_{Wilson}}\Bigr),
\qquad
\frac{dc_{w,d}}{dt} = \lambda\left.\frac{dc_{w,d}}{dt}\right|_{\mathrm{Wilson}},
\qquad
\frac{dH_d}{dt} = \lambda\left.\frac{dH_d}{dt}\right|_{\mathrm{Wilson}}
\label{eq:massdes_scale}
\end{equation}
If $\dot{m}_{Wilson} \approx 0$ but $\dot{m}_{vac} > 0$, thickness is held fixed and concentration depletes directly:
$dc_{w,d}/dt = -\dot{m}_{vac}/(MW_w H_d)$, $dH_d/dt = 0$.
Desorption rates are enforced non-positive ($dc_{w,d}/dt \le 0$, $dH_d/dt \le 0$), and the natural vacuum flux returned to the energy balance is
\begin{equation}
\dot{m}_{des,nat} = \dot{m}_{Wilson}\,\lambda = \min(\dot{m}_{Wilson}, \dot{m}_{vac}).
\label{eq:mdesnat}
\end{equation}

\subsection{MOF (dual-site Langmuir placeholder)}

Equilibrium loading $q(a_w)$ from two Langmuir sites; adsorption flux uses the Wilson Eq.~5 analog with $g_{conv}$ and $RH_{amb} - a_w(q)$. Natural desorption:
\begin{equation}
\dot{m}_{des,nat} = C_{vac}\,\bigl[P_{sat}(T_d) - P_{cond}\bigr],
\qquad
\frac{dq_d}{dt}\Big|_{\mathrm{nat}} = -\frac{\dot{m}_{des,nat}}{m_{ads}}
\end{equation}
(with inventory cap as in Eq.~\eqref{eq:massdes}). Equal mass transfer (Eqs.~\eqref{eq:equal_mass}--\eqref{eq:equal_scale}) then sets $\dot{m}_{ads}=\dot{m}_{des}$ and scales $dq_a/dt$, $dq_d/dt$.

\section{Equal mass transfer, half-cycle termination, and feedback control}

\subsection{Instantaneous equal mass transfer}

Each integration step first computes natural bed fluxes
($\dot{m}_{ads,nat}$ from the adsorption block, Eqs.~\eqref{eq:massads};
$\dot{m}_{des,nat}$ from Eqs.~\eqref{eq:massdes}--\eqref{eq:mdesnat} on the desorbing bed),
then equalizes them:
\begin{equation}
\dot{m}_{eq} = \min(\dot{m}_{ads,nat},\, \dot{m}_{des,nat}),
\qquad
\dot{m}_{ads} = \dot{m}_{des} = \dot{m}_{eq}
\label{eq:equal_mass}
\end{equation}
Loading and thickness rates on each bed are scaled by the corresponding ratio
($s_{ads} = \dot{m}_{eq}/\dot{m}_{ads,nat}$, $s_{des} = \dot{m}_{eq}/\dot{m}_{des,nat}$):
\begin{equation}
\frac{dc_{w,a}}{dt} \leftarrow s_{ads}\,\frac{dc_{w,a}}{dt}\Big|_{\mathrm{nat}},
\quad
\frac{dH_a}{dt} \leftarrow s_{ads}\,\frac{dH_a}{dt}\Big|_{\mathrm{nat}},
\quad
\frac{dc_{w,d}}{dt} \leftarrow s_{des}\,\frac{dc_{w,d}}{dt}\Big|_{\mathrm{nat}},
\quad
\frac{dH_d}{dt} \leftarrow s_{des}\,\frac{dH_d}{dt}\Big|_{\mathrm{nat}}
\label{eq:equal_scale}
\end{equation}
(and analogously $dq_a/dt$, $dq_d/dt$ for MOF).
This enforces $\dot{m}_{ads} = \dot{m}_{des}$ at every time step; integrated over a half-cycle,
\begin{equation}
\int_0^{\tau_{1/2}} \dot{m}_{ads}\,dt = \int_0^{\tau_{1/2}} \dot{m}_{des}\,dt
\label{eq:cycle}
\end{equation}
to within numerical tolerance.

\subsection{Half-cycle termination}

A half-cycle ends when the vapor-gap RH outside the desorber falls to the switch threshold,
\begin{equation}
c_r(T_d, T_{cond}) \le RH_{switch},
\label{eq:rh_switch}
\end{equation}
or when the maximum duration $\tau_{1/2,\max}$ is reached (whichever comes first).
Roles then swap: the bed that adsorbed desorbs in the next half-cycle.

\subsection{Feedback control}

There is no HTF pump to control in this variant (no actuator drives temperature matching between beds): the only feedback control is the vacuum conductance, which tracks the adsorption-limited desorption capacity (evaluated from natural, unequalized fluxes):
\begin{equation}
C_{vac} = \mathrm{clip}\!\Bigl(
C_{vac,0}\,\max\!\Bigl(0.1,\; \frac{\dot{m}_{ads,nat}}{\max(\dot{m}_{des,nat}, \epsilon)}\Bigr),\;
C_{vac,\min},\; C_{vac,\max}\Bigr)
\label{eq:ctrl_cvac}
\end{equation}
$T_a$ and $T_d$ are otherwise free-running outcomes of Eqs.~\eqref{eq:desorbing}--\eqref{eq:adsorbing}, not a controlled tracking target.

\section{Constraints}

\begin{itemize}
\item Hydrogel: $H \ge H_0$; $c_w \ge c_{w,\min}$; desorption rates non-positive ($\dot{c}_w \le 0$, $\dot{H} \le 0$); $\dot{m}_{ads} = \dot{m}_{des}$ each step (Eq.~\eqref{eq:equal_mass}).
\item MOF: $q \in [q_{\min}, q_{\max}]$.
\item Vacuum conductance clipped to physical pump limits.
\item Roles swap after each half-cycle: loading, thickness, and temperature states exchange between beds.
\item Thermal efficiency: $\eta_{th} = \dot{m}_{water}\,h_{fg} / \int \dot{Q}_{wh\to d}\,dt$ over daily operation.
\end{itemize}

\section{Default parameters (data-center baseline)}

\begin{table}[h]
\centering
\small
\begin{tabular}{@{}lll@{}}
\toprule
\textbf{Parameter} & \textbf{Symbol} & \textbf{Default} \\
\midrule
Max half-cycle duration & $\tau_{1/2,\max}$ & \SI{21600}{\second} (\SI{360}{\minute}) \\
Desorber RH switch & $RH_{switch}$ & 0.35 \\
Max depletion time & $\tau_{dep,\max}$ & \SI{600}{\second} \\
Condenser pressure & $P_{cond}$ & \SI{30}{\milli\bar} \\
Vacuum gap & --- & \SI{40}{\milli\metre} \\
Contactor thermal mass & $(\rho c_p V)_d$ & \SI{1.5e5}{\joule\per\metre\squared\kelvin} \\
Contactor-side UA & $UA_d$ & \SI{800}{\watt\per\kelvin\per\metre\squared} \\
Waste-heat inlet & $T_{wh,in}$ & \SI{58}{\celsius} \\
Waste-heat flow & $\dot{m}_{wh}$ & \SI{0.15}{\kilogram\per\metre\squared\per\second} \\
Waste-heat HX conductance & $UA_{wh}$ & \SI{1200}{\watt\per\kelvin\per\metre\squared} \\
Direct waste-heat-to-desorber equivalent UA & $UA_{wh,d}$ & \SI{480}{\watt\per\kelvin\per\metre\squared} (Eq.~\eqref{eq:uaeq}) \\
Process air & $T_{amb}$, RH & \SI{32}{\celsius}, 45\% \\
External convection & $h_{amb}$ & \SI{10}{\watt\per\metre\squared\kelvin} \\
Adsorbing-bed fin area ratio & $\phi_{fin}$ & 7.1 \\
Vacuum conductance base & $C_{vac,0}$ & \SI{8e-9}{\kilogram\per\second\per\pascal\per\metre\squared} \\
Hydrogel $H_0$ & --- & \SI{4}{\milli\metre} \\
Chamber $g$ (abs.) & $g_{chamber}$ & \SI{0.0085}{\metre\per\second} \\
Condensation enthalpy & $h_{fg}$ & \SI{2.256e6}{\joule\per\kilogram} \\
Fin area ratio (condenser) & $A_r$ & 7.1 \\
\bottomrule
\end{tabular}
\end{table}

\section*{Summary: mapping from \texttt{waste\_heat\_cycle\_lumped} (looped, 2-bed)}

\begin{table}[h]
\centering
\small
\begin{tabular}{@{}p{0.34\textwidth}p{0.30\textwidth}p{0.30\textwidth}@{}}
\toprule
\textbf{Looped variant} & \textbf{This system (no loop)} & \textbf{What changes} \\
\midrule
Loop-fluid state $T_f$ & Removed & One fewer ODE state \\
$\dot{Q}_{f\to d}$ (loop $\to$ desorber) & $\dot{Q}_{wh\to d}$ (Eq.~\eqref{eq:Qwhd}) & Waste heat couples directly through $UA_{wh,d}$ \\
$\dot{Q}_{a\to f}$ (adsorber $\to$ loop, heat recovery) & $h_{fin}A_a(T_a - T_{amb})$ & Heat of adsorption rejected to ambient, not recovered \\
$\dot{Q}_{wh\to f}$ (waste heat $\to$ loop) & Folded into $UA_{wh,d}$ (Eq.~\eqref{eq:uaeq}) & Series-combined with the old loop-to-desorber UA \\
HTF pump control ($\dot{m}_f$ tracks $T_a - T_d$) & Removed & No actuator; $T_a$, $T_d$ free-run \\
Vacuum conductance control $C_{vac}$ & Unchanged (Eq.~\eqref{eq:ctrl_cvac}) & --- \\
Sorbent mass transport (Eqs.~\eqref{eq:massads}--\eqref{eq:mdesnat}) & Unchanged & --- \\
\bottomrule
\end{tabular}
\end{table}

\section*{Integration}

Half-cycles are integrated with SciPy \texttt{solve\_ivp} (Radau). A full cycle runs two half-cycles with role swap; daily operation repeats cycles over a waste-heat / process-air profile.

\end{document}
